\chapter{Seborrheic Keratosis}

\section*{Definition}
Seborrheic keratosis (SK) is a benign, non-invasive epidermal tumor composed of proliferating keratinocytes, presenting as a well-circumscribed, waxy, stuck-on plaque with a verrucous or cerebriform surface. It is the most common benign skin tumor in adults, affecting nearly everyone over age 50.

\section*{Pathophysiology}
SKs arise from clonal proliferation of basaloid keratinocytes driven by activating mutations in FGFR3, PIK3CA, and KRAS genes in approximately 50\% of lesions. Accumulation of immature keratinocytes and pseudohorn cysts produces the characteristic keratotic architecture. The precise trigger is unknown, but UV exposure is not a major etiologic factor; hormonal influences may play a role given the increased frequency during pregnancy and estrogen therapy.

\section*{Causes}
\begin{itemize}
  \item \textbf{Aging:} Incidence increases dramatically after age 40; universal by age 70--80
  \item \textbf{Genetic:} Autosomal dominant inheritance in familial SK; FGFR3, PIK3CA, KRAS, EGFR mutations
  \item \textbf{Hormonal:} Increased during pregnancy and exogenous estrogen exposure
  \item \textbf{Inflammatory:} Leser--Tr\'elat sign --- eruptive, pruritic SKs associated with underlying malignancy (gastric adenocarcinoma, colon, breast, lymphoma)
  \item \textbf{Subtypes:} Acanthotic, hyperkeratotic, adenoid (reticulated), irritated, clonal, melanoacanthoma, stucco keratosis
\end{itemize}

\section*{When to See a Doctor}
See your doctor if:
\begin{itemize}
  \item Sudden eruption of multiple, pruritic SKs (Leser--Tr\'elat sign --- must exclude internal malignancy)
  \item Lesion changes color, bleeds, ulcerates, or becomes painful (rule out squamous cell carcinoma or melanoma)
  \item Cosmetic concern or irritation from friction (clothing, waistband, bra line)
  \item Diagnostic uncertainty: lesion with atypical features concerning for malignancy
\end{itemize}

\section*{Related Symptoms}
\begin{itemize}
  \item Actinic keratosis, viral wart, compound nevus, melanoma (especially pigmented SK mimicking melanoma)
\end{itemize}
\textbf{Important:} The stuck-on appearance and pseudohorn cysts on dermoscopy are highly characteristic of SK; however, pigmented SK may mimic melanoma and requires careful dermoscopic evaluation.

\section*{Self-Care / Home Management}
\begin{itemize}
  \item No treatment necessary for asymptomatic, benign-appearing SKs
  \item Avoid picking, scratching, or attempting to scrape off lesions, which may cause bleeding and secondary infection
  \item Sun protection does not prevent SK formation, though it may reduce actinic keratoses
\end{itemize}

\section*{How Your Doctor Will Evaluate You}
\begin{itemize}
  \item Clinical diagnosis: Sharply demarcated, waxy, stuck-on plaque with surface crypts, milia-like cysts, and comedo-like openings
  \item Dermoscopy: Comedo-like openings (crypts), milia-like cysts, fissures and ridges (brain-like pattern), hairpin vessels
  \item Skin biopsy for histopathology if diagnosis uncertain; features include acanthosis, papillomatosis, thickening of the outer skin, pseudohorn cysts
\end{itemize}

\section*{Treatment Options}
\begin{itemize}
  \item No treatment indicated for benign, asymptomatic lesions; reassurance is sufficient
  \item Cryotherapy: Liquid nitrogen with 5--10 second freeze cycle, one or two sessions for typical SK
  \item Curettage: Sharp curette after local anesthesia for larger or thicker lesions
  \item Electrocautery, laser ablation, or topical hydrogen peroxide 40\% (FDA-approved) for select lesions
  \item Histologic confirmation if any atypical features exist before destructive treatment
\end{itemize}

\clearpage
